\chapter{Lesson 8: Seeing a Doctor}

\section{Vocabulary}

\begin{center}
\begin{tabular}{llll}
\toprule
Characters & Pinyin & Part of speech & Meaning \\
\midrule
\zh{看病} & kàn bìng & verb-object & to see a doctor; (doctor) to treat \\
\zh{生病} & shēng bìng & verb-object & to get sick \\
\zh{病人} & bìngrén & noun & patient \\
\zh{医院} & yīyuàn & noun & hospital \\
\zh{疼} & téng & adjective & painful, to ache \\
\zh{头} & tóu & noun & head \\
\zh{肚子} & dùzi & noun & belly \\
\zh{嗓子} & sǎngzi & noun & throat \\
\zh{发烧} & fā shāo & verb-object & to have a fever \\
\zh{感冒} & gǎnmào & verb / noun & to catch a cold; a cold \\
\zh{咳嗽} & késou & verb & to cough \\
\zh{过敏} & guòmǐn & verb & to be allergic \\
\zh{检查} & jiǎnchá & verb & to examine \\
\zh{药} & yào & noun & medicine \\
\zh{吃药} & chī yào & verb-object & to take medicine \\
\zh{打针} & dǎ zhēn & verb-object & to get an injection \\
\zh{片} & piàn & measure word & tablet \\
\zh{次} & cì & measure word & time(s), occurrence \\
\zh{休息} & xiūxi & verb & to rest \\
\zh{好好儿} & hǎohāor & adverb & properly, well \\
\zh{健康} & jiànkāng & adjective / noun & healthy; health \\
\zh{保险} & bǎoxiǎn & noun & insurance \\
\zh{把} & bǎ & preposition & (disposal marker) \\
\zh{越来越} & yuè lái yuè & adverb & more and more \\
\zh{要不然} & yàobùrán & conjunction & otherwise \\
\zh{一\ldots{}就} & yī \ldots jiù & & as soon as \\
\zh{死了} & sǐ le & & extremely (colloquial) \\
\bottomrule
\end{tabular}
\end{center}

\section{Describing symptoms}

\begin{itemize}
  \item \zh{你怎么了？哪儿不舒服？} What's wrong? Where don't you feel well?
  \item \zh{我头疼，嗓子也疼。} I have a headache and a sore throat.
  \item \zh{我肚子疼死了。} My stomach hurts terribly.
  \item \zh{我发烧了，三十八度五。} I have a fever, 38.5 degrees.
  \item \zh{我对花生过敏。} I'm allergic to peanuts. (\zh{对} X \zh{过敏})
\end{itemize}

\section{Grammar}

\subsection{\zh{把} sentences}

The \zh{把} construction moves the object in front of the verb in order to say what
was \emph{done to} it (\zh{处置} chǔzhì ``disposal''). It is used constantly in
instructions: doctors' orders, recipes, directions.

\begin{center}
\begin{tikzpicture}[w/.style={draw, rounded corners, minimum height=0.85cm, align=center, font=\small}]
  \node[w, fill=blue!10] (s1) at (0,1.8) {\zh{你}};
  \node[w, fill=orange!20] (v1) at (1.5,1.8) {\zh{吃}};
  \node[w, fill=yellow!25] (o1) at (3.3,1.8) {\zh{这些药}};
  \node[font=\small, anchor=west] at (5.4,1.8) {S V O: ``you eat this medicine''};
  \node[w, fill=blue!10] (s2) at (0,0) {\zh{你}};
  \node[w, fill=red!15] (ba) at (1.5,0) {\zh{把}};
  \node[w, fill=yellow!25] (o2) at (3.3,0) {\zh{这些药}};
  \node[w, fill=orange!20] (v2) at (5.1,0) {\zh{吃}};
  \node[w, fill=green!15] (c2) at (6.5,0) {\zh{了}};
  \node[font=\small, anchor=west] at (7.3,0) {``take (all) this medicine''};
  \draw[->, thick, gray] (o1.south) to[bend right=15] (o2.north);
  \node[font=\scriptsize, below=0.1cm of o2] {definite object};
  \node[font=\scriptsize, below=0.1cm of c2, align=center] {something\\ after V};
\end{tikzpicture}
\end{center}

\begin{definition}[\zh{把} sentence]
Subject + \zh{把} + Object + Verb + \emph{other element}. The other element may be
\zh{了}, a resultative or directional complement, a place phrase with \zh{在/到}, a
reduplicated verb, or an indirect object with \zh{给}.
\end{definition}

\begin{itemize}
  \item \zh{请把药吃了。} Please take the medicine.
  \item \zh{我把作业写完了。} I've finished my homework.
  \item \zh{请把门关上。} Please close the door.
  \item \zh{把书放在桌子上。} Put the book on the table.
  \item \zh{你把这张表填一下。} Please fill in this form.
  \item \zh{他把钱给医生了。} He gave the money to the doctor.
\end{itemize}

\begin{important}[\zh{把} rules]
\begin{enumerate}
  \item The object is definite (specific, known): \zh{把这本书}, not *\zh{把一本书}.
  \item The verb cannot stand alone: *\zh{我把药吃} is wrong; add \zh{了} or a
        complement.
  \item Verbs with no effect on the object cannot be used: *\zh{我把他认识了},
        *\zh{把音乐喜欢}. Also not \zh{是, 有, 在, 觉得, 知道}.
  \item Negation and modals go \emph{before} \zh{把}: \zh{我没把药吃完},
        \zh{你应该把药吃了}, \zh{别把药放在冰箱里}.
  \item With \zh{放在, 放到, 送给} type verbs, \zh{把} is practically required.
\end{enumerate}
\end{important}

\subsection{\zh{越来越} and \zh{越\ldots{}越\ldots}}

\zh{越来越} + adjective/psychological verb = ``more and more''.
\begin{itemize}
  \item \zh{他的病越来越严重了。} His illness is getting more and more serious.
  \item \zh{天气越来越冷。}
  \item \zh{我越来越喜欢中文了。}
\end{itemize}
\zh{越} A \zh{越} B = ``the more A, the more B'':
\begin{itemize}
  \item \zh{雨越下越大。} The rain is getting heavier.
  \item \zh{你越休息，病好得越快。} The more you rest, the faster you'll recover.
\end{itemize}

\begin{important}[No \zh{很/非常} with \zh{越来越}]
*\zh{越来越很冷} is wrong; \zh{越来越} already expresses degree.
\end{important}

\subsection{Frequency: \zh{次}}

Number + \zh{次} goes \emph{after} the verb; with a frequency period first:
\begin{itemize}
  \item \zh{这种药一天吃三次，一次吃两片。} Take this medicine three times a day, two
        tablets each time.
  \item \zh{我去过两次医院。} I've been to the hospital twice.
  \item \zh{每四个小时吃一次。} Take it every four hours.
\end{itemize}

\subsection{Duration with double \zh{了}}

V + \zh{了} + duration + \zh{了} means the state has lasted that long \emph{and is still
continuing}; a single \zh{了} only reports the past duration.

\begin{center}
\begin{tikzpicture}[x=1cm]
  \draw[thick, -{Stealth}] (0,0) -- (10,0) node[right, font=\small] {time};
  \draw[thick] (7,0.2) -- (7,-0.2); \node[below, font=\scriptsize] at (7,-0.25) {now};
  \draw[line width=3pt, blue!60] (1,0.6) -- (4,0.6);
  \node[above, font=\small] at (2.5,0.7) {\zh{我病了三天。}};
  \node[below, font=\scriptsize, blue!60!black] at (2.5,0.45) {sick for 3 days (over)};
  \draw[line width=3pt, red!60] (4,1.8) -- (7,1.8);
  \draw[->, thick, red!60, dashed] (7,1.8) -- (8.4,1.8);
  \node[above, font=\small] at (5.5,1.9) {\zh{我病了三天了。}};
  \node[right, font=\scriptsize, red!60!black] at (8.4,1.8) {still sick};
\end{tikzpicture}
\end{center}

\subsection{\zh{要不然} ``otherwise''}

\zh{你得好好儿休息，要不然病不会好。} You must rest properly, otherwise you won't get
better.

\section{At the clinic}

\begin{example}[A visit]\leavevmode
\begin{itemize}
  \item \textbf{Doctor:} \zh{你怎么了？}
  \item \textbf{Patient:} \zh{我昨天晚上吃了一些不新鲜的东西，今天早上肚子就开始疼，疼了一天了。}
  \item \textbf{Doctor:} \zh{我检查一下。你发烧吗？}
  \item \textbf{Patient:} \zh{有一点儿。我是不是得打针？}
  \item \textbf{Doctor:} \zh{不用打针，吃药就行了。这种药一天三次，一次两片。}
  \item \textbf{Patient:} \zh{我对青霉素过敏。}
  \item \textbf{Doctor:} \zh{这不是青霉素。你回家以后把药吃了，好好儿休息，多喝水。要是越来越疼，就马上再来。}
\end{itemize}
(\zh{新鲜} xīnxiān = fresh; \zh{开始} kāishǐ = begin; \zh{青霉素} qīngméisù = penicillin;
\zh{马上} mǎshàng = immediately.)
\end{example}

\section{Culture: health care}

\begin{aside}
\begin{itemize}
  \item Chinese patients usually go straight to a hospital outpatient department rather
        than a family doctor; you register (\zh{挂号} guà hào) first and pay per visit.
  \item Traditional Chinese medicine (\zh{中医} zhōngyī) coexists with Western medicine
        (\zh{西医}); many hospitals have both departments. TCM talks about balancing
        ``hot'' and ``cold'' in the body; ``too much heat'' (\zh{上火} shàng huǒ) is blamed
        for sore throats and pimples.
  \item The standard advice for any illness is \zh{多喝热水} ``drink more hot water'', which
        has become an internet joke.
  \item At Waterloo, international students are covered by UHIP; Canadian students by
        OHIP. In Chinese, \zh{医疗保险} yīliáo bǎoxiǎn.
\end{itemize}
Traditional: \zht{醫院} hospital, \zht{藥} medicine.
\end{aside}

\section*{Review}
\begin{itemize}
  \item Symptoms: X \zh{疼}, \zh{发烧, 感冒, 咳嗽}, \zh{对} X \zh{过敏}.
  \item S \zh{把} O V + other element; definite object; negation/modals before \zh{把}.
  \item \zh{越来越} adj; \zh{越} A \zh{越} B; no \zh{很}.
  \item V + number + \zh{次}; \zh{一天三次}.
  \item V\zh{了} + duration + \zh{了} = still continuing.
  \item \zh{要不然} = otherwise.
\end{itemize}
